% =============================================================================
% CompetitiveExam/CS502/02-computer-system-organization-questions.tex
% Competitive Exam Practice 2 — GATE (Computer Science) Pattern
% Computer-System Organization and Architecture
%
% Student-facing — NO answers. Answer key: 02-computer-system-organization-answers.tex.
% Provenance: ALL questions are Original (GATE-pattern). The repository holds no
% verified GATE OS PYQ source (no local PYQ book indexed for Operating Systems),
% so no item claims to be a quoted past-year question — see the MCQ provenance
% policy in syllabi/CS502.md and .claude/rules/verification-protocol.md.
%
% Compile with: cd CompetitiveExam/CS502 &&
%   TEXINPUTS=../../Preambles;$TEXINPUTS xelatex -interaction=nonstopmode 02-computer-system-organization-questions.tex
% =============================================================================

\documentclass[11pt]{article}
\usepackage[margin=1in]{geometry}
\input{../../Preambles/header}
\coursecode{CS502}
\renewcommand{\thesection}{2.\arabic{section}}

\title{Competitive Exam Practice 2: Computer-System Organization and Architecture\\\large GATE (Computer Science) Pattern}
\author{Auqib Hamid Lone\\[0.3em]
  \emph{Department of Computer Science \& Engineering}, \emph{GCET Ganderbal}}
\date{\today}

\begin{document}
\maketitle

\noindent Each question carries 1--2 marks. MCQ = single correct option; NAT = Numerical Answer Type. \emph{All questions in this set are original, written in the GATE pattern --- the repo holds no verifiable GATE OS past-year source.}

\begin{enumerate}
\item[\textbf{Q1}] \textit{[Original, GATE-pattern]} \textbf{[MCQ, 1 mark]} Which statement about control transfer to the kernel is \textbf{TRUE}?

\begin{enumerate}
  \item[(A)] An interrupt is synchronous with the instruction that raised it, whereas a trap is asynchronous.
  \item[(B)] A divide-by-zero fault is delivered as a hardware interrupt, because the arithmetic unit signals it.
  \item[(C)] A system call enters the kernel through a trap, whereas a disk-completion signal is delivered as an interrupt.
  \item[(D)] Both interrupts and traps can occur only at an instruction boundary chosen by the running process.
\end{enumerate}

\item[\textbf{Q2}] \textit{[Original, GATE-pattern]} \textbf{[NAT, 2 marks]} An interrupt vector table starts at physical address $4096$, and entry $i$ occupies the 8 bytes beginning at address $4096 + 8i$. The CPU receives interrupt number $40$. The physical address at which the CPU fetches this interrupt's handler address is \rule{1.5cm}{0.4pt}.

\item[\textbf{Q3}] \textit{[Original, GATE-pattern]} \textbf{[NAT, 2 marks]} A network adapter raises $4{,}000$ interrupts per second. Saving and restoring the context and dispatching to the handler costs $4.5\ \mu\text{s}$ per interrupt, and the device handler itself runs for $8\ \mu\text{s}$. Ignoring all other overhead, the percentage of CPU time consumed by this adapter's interrupts is \rule{1.5cm}{0.4pt}\,\%.

\item[\textbf{Q4}] \textit{[Original, GATE-pattern]} \textbf{[MCQ, 1 mark]} Which statement about maskable interrupts and priority is \textbf{TRUE}?

\begin{enumerate}
  \item[(A)] Clearing the CPU's interrupt-enable flag also prevents delivery of non-maskable interrupts.
  \item[(B)] A masked interrupt request is discarded, so it can never be serviced at a later time.
  \item[(C)] While a high-priority handler runs with lower-priority interrupts masked, a pending lower-priority request is normally serviced after the high-priority handler returns.
  \item[(D)] Interrupt priority is determined solely by the order in which handlers are installed in software; the interrupt controller treats all requests equally.
\end{enumerate}

\item[\textbf{Q5}] \textit{[Original, GATE-pattern]} \textbf{[NAT, 2 marks]} A magnetic disk rotates at $7{,}200$ RPM. Its average seek time is $6\ \text{ms}$, its sustained transfer rate is $128\ \text{MB/s}$, and one block is $4\ \text{KB}$. The average time to read one block, in milliseconds rounded to two decimal places, is \rule{1.5cm}{0.4pt}.

\item[\textbf{Q6}] \textit{[Original, GATE-pattern]} \textbf{[MCQ, 1 mark]} Which statement about volatility in the memory hierarchy is \textbf{TRUE}?

\begin{enumerate}
  \item[(A)] DRAM is non-volatile, which is why a program's variables survive a power failure.
  \item[(B)] SRAM caches are non-volatile, whereas DRAM is volatile.
  \item[(C)] Registers and DRAM lose their contents without power, whereas magnetic disk and NAND flash retain theirs.
  \item[(D)] NAND flash is volatile, but it has a lower cost per byte than magnetic disk.
\end{enumerate}

\item[\textbf{Q7}] \textit{[Original, GATE-pattern]} \textbf{[NAT, 2 marks]} A two-level cache hierarchy has L1 hit time $1\ \text{ns}$ and L1 hit rate $90\%$. After an L1 miss is detected, an L2 access costs $8\ \text{ns}$, and $80\%$ of L1 misses hit in L2. An access that misses in both caches costs a further $60\ \text{ns}$ for main memory. The average memory access time (AMAT), in nanoseconds, is \rule{1.5cm}{0.4pt}.

\item[\textbf{Q8}] \textit{[Original, GATE-pattern]} \textbf{[MCQ, 1 mark]} Which statement about polling, interrupt-driven I/O, and DMA is \textbf{TRUE}?

\begin{enumerate}
  \item[(A)] Polling never reads the device's status register, so it cannot determine whether the device is ready.
  \item[(B)] In interrupt-driven I/O the device signals the CPU when service is needed, so CPU cost is incurred per event rather than by busy-waiting.
  \item[(C)] DMA removes all cache-coherence concerns, because the controller writes directly into main memory.
  \item[(D)] DMA is cheaper than a single interrupt for a one-byte transfer from a slow keyboard, since setup cost is negligible.
\end{enumerate}

\item[\textbf{Q9}] \textit{[Original, GATE-pattern]} \textbf{[NAT, 2 marks]} A $2\ \text{GHz}$ CPU spends $2{,}000$ cycles programming a DMA controller and $1{,}000$ cycles in the completion interrupt handler. The controller then moves $1\ \text{MB}$ at $200\ \text{MB/s}$, a transfer lasting $5\ \text{ms}$. The DMA-related CPU time, expressed as a percentage of this $5\ \text{ms}$ transfer time and rounded to two decimal places, is \rule{1.5cm}{0.4pt}\,\%.

\item[\textbf{Q10}] \textit{[Original, GATE-pattern]} \textbf{[MCQ, 1 mark]} Which statement about multiprocessor, multicore, and cluster architectures is \textbf{TRUE}?

\begin{enumerate}
  \item[(A)] In a multicore processor, all cores share a single set of architectural registers.
  \item[(B)] SMP requires every processor to have private main memory; otherwise the processors cannot run one OS image.
  \item[(C)] A cluster connects independent computers over a network and can continue service after one whole machine fails.
  \item[(D)] A cluster is a single chip containing several cores that share a last-level cache.
\end{enumerate}

\end{enumerate}

\end{document}
